\chapter*{Introduzione}

\addcontentsline{toc}{chapter}{Introduzione}

I capitoli dell'elaborato coincidono con gli esercizi svolti. Ogni capitolo è suddiviso in paragrafi, a cui corrispondono i singoli punti degli esercizi (Es. 1.1, 1.2), dove prima viene esaminata l'architettura del sistema, poi il codice dei componenti e infine i risultati di simulazione o sintesi su scheda. La \textit{board} in dotazione al gruppo è la Nexys A7-100T che ha un \textit{clock} di base di 100MHz.
Inoltre, sono disponibili, come materiale multimediale aggiuntivo, dei video dimostrativi relativi alle implementazioni su board degli esercizi. Il nome dei file corrisponde al numero dell'esercizio, e sono consultabili al seguente link:

\url{https://drive.google.com/drive/folders/1lWYjplbmSitBj_Bi_736sA2Rr00uTwj_?usp=share_link}

